\chapter{Оркестратор и цикл задачи}

\section{Роль и её границы}

Оркестратор планирует: декомпозирует цели в задачи с контрактами, назначает продьюсеров, инициирует ревью и интеграцию. Это прикладная роль, а не роль в консенсусе. Все её действия суть записи, проходящие правила валидаторов. Полезно перечислить, чего оркестратор не может:

\begin{itemize}
\item не имеет доступа к среде помимо записей: у него нет shell, среда исполняет команды воркеров;
\item не выбирает ревьювера: выбор вычисляется из случайности, полученной после того, как подбирать уже нечего;
\item не интегрирует без одобрения и проверок, адресованных к кандидату;
\item не выходит за capability и не меняет условия приёмки иначе как новой версией контракта с авторизацией;
\item не может заблокировать остановку.
\end{itemize}

\section{Аренда и fencing}

Роль оформлена записью аренды: нода, эпоха, случайный fencing-токен, предшествующая аренда, причина (создание, плановая передача, принудительный захват, восстановление), чекпоинт как база передачи и область. Аренда продлевается heartbeat. Каждая запись оркестратора несёт эпоху и токен; валидаторы отвергают записи с устаревшей арендой, но сохраняют их как ветку расхождения для аудита (инвариант I7).

\begin{exbox}[Уснувший ноутбук]
Оркестратор на ноутбуке уходит в сон. Heartbeat пропадает, аренда истекает, роль по ротации переходит следующей кор-ноде с новой эпохой и новым токеном. Ноутбук просыпается и отправляет назначение с прежним токеном. Валидаторы его не подписывают, но записывают в ветку расхождений. Среда, даже если запись каким-то образом до неё дошла бы, не исполнит её без сертификата кворума и без проверки актуальности эпохи на границе исполнения.
\end{exbox}

\section{Смена эпохи}

Смена происходит по таймауту heartbeat, по исчерпанию ёмкости оркестратора, по решению кворума или по обнаружению нарушения правил. Новый оркестратор определяется детерминированной ротацией; византийская нода пропускается следующей сменой. Раз в $k$ эпох смена происходит холодно независимо от причины.

\term{Плановая передача} требует доказательства покоя: ввод заблокирован, нет незавершённых вызовов, дочерних процессов и незавершённых операций раннера, дайджест хранилища стабилен. После этого сессия передаётся преемнику, при необходимости с конвертацией между форматами агентов, и передача записывается в журнал с хешем снимка.

\term{Принудительный захват} есть отдельное явное действие, а не плановая передача, которая не удалась. Сессия не передаётся, предшественник отсекается токеном, преемник стартует из журнала.

\begin{whybox}[Почему при захвате сессия не передаётся]
Уходящий оркестратор может быть тем самым византийским узлом. Переданная им сессия может быть отравлена, и преемник, унаследовавший такой контекст, становится византийским сам: компрометация переезжает вместе с ролью и обесценивает арифметику $f$. Поэтому враждебная смена всегда холодная, а при плановой передаче переданная сессия недоверена и расхождение с журналом разрешается в пользу журнала.
\end{whybox}

\section{Состояние и чекпоинты}

Журнал содержит план, контракты, назначения и результаты, и этого достаточно для вступления преемника. Сессия модели есть кэш. Оркестратор периодически сохраняет снимок сессии в хранилище с сертификатом доступности и фиксирует его хеш записью чекпоинта; после сбоя преемник загружает снимок и докатывает записи журнала после него. Дисциплина для оркестратора: решения и обоснования пишутся в журнал структурно, по ходу работы; конвертация сессий между агентами теряет часть контекста, и только структурная запись делает сессию по-настоящему одноразовой. Эксперимент, проверяющий это напрямую, описан в главе 11: остановить оркестратора, удалить его сессию, проверить, продолжит ли преемник.

\section{Несколько оркестраторов}

Аренда несёт область, поэтому одновременно могут действовать несколько аренд с дизъюнктными областями. Задача, затрагивающая пересечение областей, принадлежит аренде, покрывающей объединение, либо создаётся кор-нодой напрямую.

Когда один борд процесса ведут несколько операторов, 28 сентября 2026 года решено направление: борд остаётся один на процесс, и у него или у совместной работы один \term{принимающий писатель} с fencing по эпохе. Каждая типизированная операция несёт актора, борд, версию конфигурации, область и эпоху аренды, ожидаемую ревизию и идентификатор операции. Эпоха проверяется при запуске агента, при финальной записи и при интеграции; писатель, потерявший полномочия, приостанавливает новые записи. Это та же аренда с fencing-токеном, применённая к одному борду. Решение идёт в трек waggle CM2 вместе с арендами областей (waggle, \texttt{spec/waggle.md} v5.2, §12), а назначается ли писатель по владению процессом или выбирается участниками, пока открыто. В терминах главного цикла узла оркестратор процесса есть его корневой оркестратор, а оркестраторы на других узлах работают арендованными поддеревьями и второго корневого не заводят (curator-keeper, \texttt{spec/keeper.md} v0.5.1, §3.2, §4).

\section{Цикл задачи}

\begin{center}
\begin{tikzpicture}[node distance=4mm,scale=0.9,every node/.style={transform shape}]
\node[box] (s1) {Авторизована};
\node[box,right=of s1] (s2) {Назначена};
\node[box,right=of s2] (s3) {Ревизия $R$};
\node[box,right=of s3] (s4) {Обратная связь};
\node[box,below=13mm of s1] (s5) {Кандидат $(M, I)$};
\node[box,right=of s5] (s6) {Ревью $I$};
\node[box,right=of s6] (s7) {Проверки $I$};
\node[box,right=of s7] (s8) {Интеграция};
\node[box,right=of s8] (s9) {Вестинг};
\draw[arr] (s1) -- (s2); \draw[arr] (s2) -- (s3); \draw[arr] (s3) -- (s4);
\draw[arr] (s4.south) -- ++(0,-5mm) -| (s5.north);
\draw[arr] (s5) -- (s6); \draw[arr] (s6) -- (s7); \draw[arr] (s7) -- (s8); \draw[arr] (s8) -- (s9);
\draw[arr] (s4.north) -- ++(0,5mm) -| node[lbl,above,pos=0.25] {не более $N$ раундов} (s3.north);
\draw[arr] (s7.south) -- ++(0,-4mm) -| node[lbl,below,pos=0.3] {отклонено или голова ушла вперёд: пересборка кандидата} (s5.south);
\end{tikzpicture}
\end{center}

\subsection{Авторизация и назначение}

Оркестратор декомпозирует цель в задачи с контрактами; задача совершается записью с сертификатом кворума. Рабочие ноды сообщают в heartbeat остаток ёмкости и окно его сброса. Оркестратор назначает продьюсера с учётом ёмкости, свободного слота и удерживаемых областей. Среда открывает санбокс по записи назначения.

\subsection{Ревизия и обратная связь}

Продьюсер работает в санбоксе и представляет хеш коммита $R$ над базой $B$. Подпись ревизии ставится на машине продьюсера, ключ никогда не покидает ноду (инвариант I6). Голова цепочки receipts входит в запись.

Ранняя обратная связь ревьювера на $R$ допустима и полезна: продьюсер получает находки и представляет новую ревизию, число оплачиваемых раундов ограничено $N$. Но обратная связь не есть одобрение. Одобрение относится только к кандидату.

\subsection{Интеграционный кандидат}

\begin{center}
\begin{tikzpicture}[node distance=9mm]
\node[box] (b) {$B$\\[1pt]\footnotesize база задачи};
\node[box,right=of b] (r) {$R$\\[1pt]\footnotesize ревизия продьюсера};
\node[box,below=8mm of b] (m) {$M$\\[1pt]\footnotesize текущая голова main};
\node[box,right=of m] (i) {$I$\\[1pt]\footnotesize кандидат = слияние $M$ и $R$};
\draw[arr] (b) -- (r);
\draw[arr] (b) -- node[lbl,left] {main ушла вперёд} (m);
\draw[arr] (m) -- (i);
\draw[arr] (r) -- (i);
\node[lbl,right=4mm of i,align=left,text width=52mm] {одобрение и проверки адресованы к $I$, версии контракта, политике и окружению; после сдвига головы кандидат перестраивается, ревью и проверки повторяются};
\end{tikzpicture}
\end{center}

Среда строит неизменяемого кандидата $I$ над текущей головой $M$ и записывает пару $(M, I)$. Множество кандидатов в ревьюверы замораживается до появления хеша $I$; случайность получается после, из порогового маяка кор-нод или проверяемой случайной функции; ревьювер вычисляется из идентификатора задачи, хеша $I$ и случайности, и любой валидатор пересчитывает выбор.

\begin{whybox}[Почему одобрение нельзя переносить с ревизии на кандидата]
Во второй редакции одобрение $R$ переносилось на $I$, если совпадал patch-id дифа. Это было ошибкой. Одинаковый патч на другой базе может иметь другой смысл: одна задача изменила контракт функции, другая добавила её вызов, патч второй не изменился, но предположение, на котором стояло ревью, уже неверно. Исполняемая модель фрагмента интеграции (глава 11) находит ровно этот сценарий как нарушение, когда правило привязки к кандидату отключено. Поэтому финальное одобрение относится к неизменяемому $I$. Для людей повторное ревью после каждого перебазирования было бы неподъёмным, и очереди интеграции его не делают; для моделей оно стоит токены, и правило посильно. Перенос допустим только как расширение вне гарантии, при детерминированном условии: дельта базы $B..M$ не пересекается ни с write scope задачи, ни с множеством путей, которые продьюсер читал по receipts. Проверки на $I$ повторяются всегда.
\end{whybox}

\subsection{Ревью и проверки кандидата}

Ревьювер получает диф $M..I$ и подписывает одобрение, адресованное к $I$, версии контракта, политике и окружению. Проверки из профиля контракта выполняются на $I$ по инициативе ноды, третьей по отношению к продьюсеру и ревьюверу; критические наборы повторяет верификатор, если он есть. Каждая аттестация имеет один из трёх исходов; принятие требует исхода пройдена по всем свойствам профиля. При отклонении кандидата продьюсер получает находки, и цикл возвращается к ревизии.

\subsection{Интеграция}

Оркестратор вносит запись об интеграции с парой $(M, I)$. Валидаторы проверяют правило 5. Среда, получив сертификат кворума, атомарно заменяет голову основной ветки с $M$ на $I$. Если голова уже не равна $M$, сертификат неприменим: кандидат перестраивается, ревью и проверки повторяются (инвариант I3). Порядок интеграции задаёт очередь.

\subsection{Вестинг, превышение $N$, замена}

Начисление за задачу закрепляется, когда $I$ прожил $k$ эпох без отката. При превышении $N$ раундов оркестратор обязан выбрать из фиксированного меню: передекомпозировать задачу, сменить продьюсера, сменить ревьювера или эскалировать человеку; выбор записывается. Замена продьюсера происходит при потере ноды, исчерпании ёмкости или по решению из меню; преемник получает сессию предшественника с правилом недоверия к ней.

\section{Внешние эффекты}

Операции раннера, подпись, публикация, устройства, данные, имеют эффекты вне среды, и журнал авторизаций сам по себе не определяет их семантику при сбоях. Сценарий: публикация авторизована, раннер её выполняет и падает до сохранения результата; новый оркестратор видит авторизацию без завершения. Можно ли повторить?

\begin{itemize}
\item Каждая операция имеет идентификатор и ключ идемпотентности. Повтор ссылается на тот же идентификатор и параметры; для получателей, поддерживающих идемпотентность, повтор не создаёт дополнительного эффекта.
\item Исход записывается как выполнена, отклонена или не установлен. Сбой раннера до записи исхода даёт не установлен.
\item При исходе не установлен слепой повтор запрещён: требуется сверка фактического состояния получателя либо разрешение человеком, и результат сверки записывается.
\item Актуальность разрешения проверяется на границе исполнения. Операция, авторизованная в прошлой эпохе или до остановки, не исполняется без новой авторизации. Уже запущенная операция при смене аренды доводится до записи исхода; при остановке прерывается там, где это безопасно, иначе доводится и записывается.
\end{itemize}

\begin{warnbox}
Старая корректная авторизация не равна действующему разрешению на любое будущее исполнение. Правило актуальности проверяется демоном среды и раннером в момент исполнения, а не только валидаторами в момент записи.
\end{warnbox}

\section{Одна задача, все записи}

\begin{exbox}[Задача из сценария, от цели до вестинга]
\begin{enumerate}
\item Цель: перенос модуля синхронизации; область \texttt{Sources/Sync/**}; класс обычный. Подпись кор-ноды, QC.
\item Capability оркестратору: та же область, не больше шести одновременных задач.
\item Задача \texttt{TASK-7fk2}: контракт версии 1 (write scope \texttt{Sources/Sync/Transport/**}, критерии, профиль: сборка, регрессия, ревью), оценка 5, digest плана.
\item Назначение: нода контрибьютора А, роль продьюсера, аренда эпохи 41, слот 7. Среда открывает санбокс.
\item Ревизия $R$ над базой $B$, подпись А, голова цепочки receipts.
\item Обратная связь: нода В, две находки. Ревизия $R'$ от А.
\item Кандидат: голова $M$, слияние $I$; записана пара $(M, I)$.
\item Заморозка кандидатов в ревьюверы: ноды Б, В, Г. Случайность эпохи 41 после хеша $I$.
\item Ревью $I$: нода В, одобрено; адресовано к $I$, контракту версии 1, политике, окружению.
\item Аттестации: нода Г инициировала сборку и регрессию на $I$, исход пройдена; верификатор подтвердил критический набор.
\item Интеграция: пара $(M, I)$, QC. Среда: голова была $M$, стала $I$.
\item Вестинг через $k$ эпох; аттестации репутации А, В, Г.
\end{enumerate}
Двенадцать записей, из них ни одной, которую можно было бы подделать одной подписью на L1 и выше, и ни одной, для которой валидатору понадобилась бы модель.
\end{exbox}
